\cleardoublepage

\section{系统设计与实现}

% ============================================================
% 本章约15000-18000字，是毕业设计"本人工作"的核心部分
% 后端各模块各一节，前端总体一节
% 技术原理在各模块实现中顺带介绍，不再单独成章
% ============================================================

% -----------------------------------------------------------
% §3.1 系统总体架构设计
% -----------------------------------------------------------

\subsection{系统总体架构设计}

\subsubsection{技术选型}

\par 本项目的后端技术栈围绕 Spring 生态进行选型。后端核心框架选用 Spring Boot 3.5.7，提供自动配置、嵌入式容器和起步依赖管理能力。开发语言采用 Java 21，其虚拟线程特性对 AI 流式响应场景下的高并发处理具有显著优势。数据持久化选用 PostgreSQL 15，该版本在 JSONB 数据类型和并行查询方面表现优异；缓存层选用 Redis 7，承担会话状态缓存、令牌黑名单和 AI 检查点持久化职责。ORM 框架采用 MyBatis Plus 3.5.14，封装了通用 CRUD 和分页能力。AI 能力接入选用 Spring AI Alibaba 1.1.2.0，提供统一的聊天模型抽象和基于图结构的工作流编排能力，可便捷接入通义千问和 DeepSeek 等大语言模型。对象映射层选用 MapStruct 1.6.3，在编译期生成类型安全的映射代码，避免运行期反射的性能损耗。其他关键组件包括 jjwt 令牌库、libphonenumber 号码解析、uuid-creator 时间有序 UUID 生成、Flyway 数据库版本管理和 SpringDoc 接口文档自动生成。项目部署采用 Docker 容器化方案，通过 Docker Compose 编排各服务组件。

\subsubsection{分层架构设计}

\par 系统整体采用领域驱动设计的分层架构，自上而下划分为接口层、业务层、领域层、基础设施层和数据库层五个层次，每层仅向下依赖，避免跨层调用和循环依赖。图 \ref{fig:architecture} 展示了系统的总体架构。

\begin{figure}[htbp]
    \centering
    \includegraphics[width=\linewidth]{architecture.png}
    \caption{系统总体架构图}
    \label{fig:architecture}
\end{figure}

\par 接口层负责接收外部请求并进行协议转换，所有外部访问统一通过 RESTful API 接入，返回 JSON 格式的标准化响应，仅做参数校验、权限检查和响应封装。业务层承载各业务领域的用例实现和领域规则编排，模块内部通过明确定义的接口协作。领域层通过 Repository 接口抽象持久化细节，领域实体封装业务规则和数据约束。

\par 基础设施层实现持久化接口，包含基于 MyBatis Plus 的数据访问、基于 Spring AI Alibaba 的 AI 引擎接入以及阿里云短信和语音服务的客户端封装。数据库层采用双存储策略，PostgreSQL 作为主存储保证数据一致性，Redis 作为缓存和会话存储，两层之间通过应用程序协调读写。

\subsubsection{模块划分}

\par 依据领域驱动设计的限界上下文划分原则，系统按业务领域拆分为六个独立模块，各模块拥有独立的包结构和持久化实现，通过接口进行交互。

\par 公共模块提供通用基础设施，包括 Jackson 序列化配置、统一异常体系、持久化基类、自定义校验注解等。用户模块负责手机号注册、短信验证码、JWT 认证等用户生命周期管理。项目管理模块涵盖项目、任务、标签三个核心实体及数据库视图聚合查询。智能体模块负责大语言模型接入、工作流编排和工具调用，是系统的技术核心。安全模块通过过滤器链实现认证鉴权和令牌黑名单管理。语音服务模块承担阿里云语音服务的临时令牌生成与缓存，支持前端直接连接语音识别与合成服务。

\par 各模块之间的依赖关系遵循严格的单向依赖原则。公共模块作为最底层的基础设施，不依赖任何业务模块；业务模块之间通过服务接口进行调用，禁止直接访问其他模块的持久化层。这种依赖约束确保了系统的模块化程度，当某个业务模块需要调整或替换时，影响范围被限制在该模块内部，不会波及其他模块。例如，当未来需要更换短信服务提供商时，只需修改用户模块中短信服务的客户端实现，其他模块完全不受影响。

% -----------------------------------------------------------
% §3.2 数据库设计
% -----------------------------------------------------------

\subsection{数据库设计}

\par 数据库设计采用关系型模型，遵循规范化原则在消除数据冗余与保证查询效率间取得平衡。设计围绕用户、项目、任务等核心业务实体展开，通过合理的主外键关联和索引策略确保数据完整性并优化高频查询响应性能。共包含七张核心数据表和一张聚合视图。图 \ref{fig:database} 展示了核心实体之间的关联关系。

\begin{figure}[htbp]
    \centering
    \includegraphics[width=\linewidth]{database.png}
    \caption{数据库实体关系图}
    \label{fig:database}
\end{figure}

\par 用户表是系统的根实体，存储用户的基本身份信息，包括角色、手机号、昵称、头像和时区标识。手机号字段设置唯一约束，作为用户的登录凭证。时区标识用于前端将 UTC 时间戳转换为用户本地时间。项目表记录用户创建的每个项目，包含名称、描述、优先级、状态、复杂度和截止日期等业务属性。其中优先级字段使用 SMALLINT 类型，数值越小表示优先级越高，便于排序查询。复杂度字段用于智能体模块判断项目处理策略，分为简单和复杂两个级别。任务表以项目为维度进行任务拆分，每条记录关联到一个具体项目，支持状态流转和批量更新。标签表提供用户自定义的分类能力，每个标签包含名称和颜色两个属性，颜色以七字符十六进制格式存储。项目与标签之间是多对多关系，通过关联表进行桥接，该关联表同时记录用户标识，便于权限过滤。

\par 会话表和聊天消息表共同支撑智能体模块的对话功能。会话表代表一次完整的 AI 对话上下文，包含会话名称和 AI 生成的项目草稿视图。草稿视图以 JSON 文本形式存储，用于保存智能体在对话过程中为用户构建的项目结构化数据。聊天消息表记录会话中的每条消息，通过类型字段区分用户消息和 AI 回复。

\par 所有表均采用 UUID 作为主键，使用 uuid-creator 库生成第七版 UUID。该版本 UUID 基于时间排序，在保持全局唯一性的同时具备良好的插入性能，避免了传统随机 UUID 导致的数据页频繁分裂问题。表与表之间通过外键建立关联，并配置级联删除策略，当父表记录被删除时，所有关联的子表记录自动清理，保证数据一致性。时间戳字段统一使用 TIMESTAMPTZ 类型，以 UTC 格式存储，消除时区转换带来的歧义。

\par 为了优化前端的数据获取效率，设计了一个数据库视图用于聚合项目及其关联的标签和任务数据。该视图通过 JSONB 聚合函数将多行标签记录和任务记录分别聚合成 JSONB 数组，作为项目的两个派生列返回。前端通过一次查询即可获得项目的完整信息，避免了传统的 N+1 查询问题。视图的底层基表建立了复合索引，覆盖最常用的用户标识加优先级、用户标识加截止日期等查询条件，确保视图查询的响应时间在可接受范围内。

\par 数据库的版本管理采用 Flyway 工具，通过迁移脚本统一维护表结构、索引和视图的创建与变更。所有数据库对象的定义集中在一个初始化脚本中，新环境部署时 Flyway 自动按版本顺序执行脚本，保证各环境的数据库结构一致。这种版本化管理使得数据库变更以可追踪、可回滚的方式进行，避免了手动修改时常见的遗漏和冲突问题，是团队协作和持续集成流程中的重要基础设施。

\par 从数据一致性角度来看，所有表之间通过外键建立了完整的关联关系，并配置了级联删除策略。当用户被删除时，其关联的项目、任务、标签和会话记录会自动清理，避免产生数据孤儿。时间戳字段统一使用 TIMESTAMPTZ 类型以 UTC 格式存储，前端根据用户表中记录的时区标识将 UTC 时间转换为本地时间，确保不同时区用户看到的时间都符合其所在地习惯。复合索引的设计覆盖了最常用的查询场景，如按用户标识加优先级排序、按用户标识加截止日期筛选等，确保在大数据量下查询性能仍保持在可接受范围内。

% -----------------------------------------------------------
% §3.3 公共基础设施层设计
% -----------------------------------------------------------

\subsection{公共基础设施层设计}

\subsubsection{统一响应封装}

\par 系统在公共模块中定义了统一的响应结构和异常处理机制。所有接口的返回值都包装在一个泛型响应对象中，包含状态码、提示消息和实际数据三个字段。状态码采用五位数字编码，前两位表示错误类别，后三位表示具体错误，例如以四开头的状态码表示客户端请求错误，以五开头的状态码表示服务端内部错误。异常处理由全局异常处理器统一接管，该处理器通过框架提供的注解拦截控制器层抛出的所有异常，根据异常类型映射到对应的错误码和提示消息。错误码枚举按业务领域划分为不同区间，用户模块的错误码集中在四万号段，项目管理模块集中在五万号段，智能体模块集中在六万号段，通用错误码集中在三万号段。这种区间划分使得在日志分析时能够根据错误码快速定位出错的业务领域。

\subsubsection{持久化基类设计}

\par 公共模块中设计了两个持久化抽象基类，分别面向实体表操作和数据库视图操作。实体仓库基类封装了标准的增删改查方法，各业务模块只需继承基类并指定领域实体类型、持久化对象类型和数据访问映射器类型三个泛型参数，即可获得完整的持久化能力。视图仓库基类面向只读查询场景，提供查询相关方法并在内部完成持久化对象到领域实体的转换。两个基类的泛型设计使得编译器能够在编译期检查类型匹配关系，避免运行期的类型转换错误。基类还预留了钩子方法，允许子类在标准操作的前后插入自定义逻辑，例如某些业务实体在插入前需要自动生成特定的业务编码。

\subsubsection{领域实体转换体系}

\par 系统在公共模块中设计了一套分层的实体转换体系，通过两层转换器实现领域实体与持久化对象之间的双向映射。第一层是字段级转换器，负责处理值类型差异，例如 UUID 与字符串、时区对象与字符串、JSONB 数据与结构化对象之间的转换。第二层是对象级转换器，使用 MapStruct 工具在编译期生成类型安全的映射代码，开发人员只需定义转换接口并用注解标注字段映射关系，编译器在构建阶段自动生成实现类。

\par 这套转换体系的最大优势在于实现了领域模型与数据库模式的完全解耦。领域模型可以根据业务需求自由调整结构，不受数据库表结构的约束；持久化模型可以根据存储优化需要调整列类型，而不会影响业务逻辑。两个模型的变更只需在转换器中更新映射规则，对上层的业务服务和下层的存储引擎都是透明的。当未来需要更换数据库或调整表结构时，修改范围可以局限在持久化层和转换器层，业务逻辑层无需任何改动。以项目实体为例，领域模型中的 UUID 标识与数据库表中的字符串列通过字段级转换器完成映射，对象级转换器负责整体字段的复制和转换，当需要增加计算属性时只需修改领域模型和转换配置，数据库结构完全不受影响。

% -----------------------------------------------------------
% §3.4 用户模块设计与实现
% -----------------------------------------------------------

\subsection{用户模块设计与实现}

\subsubsection{用户注册与SMS验证码登录}

\par 用户模块采用手机号加短信验证码的登录方式，避免了传统密码体系带来的记忆负担和安全风险。整个登录流程分为四个步骤：用户输入手机号请求发送验证码，服务端调用短信服务发送六位验证码，用户输入验证码提交登录请求，验证通过后生成访问令牌返回。图 \ref{fig:user-pages} 展示了用户模块相关的前端页面。

\begin{figure}[htbp]
    \centering
    \includegraphics[width=0.2\linewidth]{login.png}
    \hspace{0.1\linewidth}
    
\includegraphics[width=0.2\linewidth]{home.png}
    \hspace{0.1\linewidth}
    \includegraphics[width=0.2\linewidth]{mine.png}
    \caption{用户模块前端页面}
    \label{fig:user-pages}
\end{figure}

\par 手机号格式校验采用 libphonenumber 库处理，该校验库内置了全球各国的号码格式规则，能够识别号码的国家代码和本地格式是否合法。验证码的安全机制包含有效期限制和频率限制两层保护：有效期限制确保验证码在固定时间窗口后自动失效；频率限制防止同一手机号在单位时间内高频请求验证码。短信服务集成阿里云号码认证服务，支持生产模式和测试模式切换，测试模式下返回固定验证码便于开发调试。首次登录的用户在验证通过后自动创建用户记录，实现无感注册。

\subsubsection{JWT认证实现}

\par 系统采用 JSON Web Token 实现无状态认证。访问令牌由头部、载荷和签名三部分组成，头部包含令牌类型和签名算法，载荷包含用户标识、角色和签发时间等声明，签名使用密钥对前两部分进行哈希运算确保内容未被篡改。令牌的生成和验证由认证服务负责，使用 jjwt 库采用 HMAC-SHA256 签名算法，默认有效期为二十四小时。

\par 考虑到用户主动登出的场景，系统引入了令牌黑名单机制。用户登出时当前令牌被加入缓存黑名单，此后携带该令牌的请求会被拦截拒绝。黑名单设置了与令牌剩余有效期相同的过期时间，定期自动清理已过期条目。令牌载荷预留了自定义声明扩展点，密钥通过环境变量注入，生产环境不硬编码敏感信息，所有认证接口均通过加密通道传输。

\subsubsection{MyBatis Plus与PostgreSQL持久化}

\par 用户模块的数据持久化采用 MyBatis Plus 作为对象关系映射框架，该框架封装了通用 CRUD 和分页能力，开发者只需继承基类即可获得完整的持久化能力。PostgreSQL 作为主存储数据库，支持完整的事务机制和 JSONB 数据类型，所有业务表的主键采用第七版 UUID 生成策略，该版本 UUID 基于时间排序，具备良好的插入性能。时区处理方面，所有时间戳以 UTC 格式存储，前端根据用户时区标识转换为本地时间。

\par 数据库版本管理采用 Flyway 工具，所有表定义、索引和视图集中在一个初始化脚本中，新环境部署时自动按版本顺序执行。当需要变更表结构时，开发者创建新的迁移脚本，确保变更以可追踪、可回滚的方式应用，避免了手动修改数据库时常见的遗漏和冲突问题。MyBatis Plus 的条件构造器还为复杂查询提供了类型安全的构建方式，开发者通过链式调用构造查询条件，框架自动转换为底层 SQL 语句，编译器在编译期即可检查字段名称和类型的正确性，消除了字符串拼接带来的语法错误风险。

% -----------------------------------------------------------
% §3.5 项目管理模块设计与实现
% -----------------------------------------------------------

\subsection{项目管理模块设计与实现}

\subsubsection{项目与任务管理}

\par 项目管理模块围绕项目和任务两个核心实体提供生命周期管理能力。项目具有名称、描述、优先级、状态、复杂度和截止日期等属性，优先级采用小整数类型，数值越小优先级越高，便于排序查询。状态机包含待处理、进行中、已完成和已归档四个状态，流转遵循业务规则约束。任务是项目的细粒度拆分单元，包含待办、进行中和已完成三个状态，支持批量更新和事务包裹。复杂度字段是项目管理模块与智能体模块之间的关键关联，用户在创建项目时标记复杂度级别，直接影响智能体的处理策略选择。图 \ref{fig:project-pages} 展示了项目管理相关的前端页面。

\begin{figure}[htbp]
    \centering
    \includegraphics[width=0.2\linewidth]{dashboard.png}
    \hspace{0.05\linewidth}
    
\includegraphics[width=0.2\linewidth]{project.png}
    \hspace{0.05\linewidth}
    \includegraphics[width=0.2\linewidth]{focus.png}
    \hspace{0.05\linewidth}
    \includegraphics[width=0.2\linewidth]{edit.png}
    \caption{项目管理模块前端页面}
    \label{fig:project-pages}
\end{figure}

\par 项目服务层封装了完整的生命周期逻辑，创建时进行业务规则校验，更新时将变更应用到领域实体后持久化，删除时采用级联策略清理关联记录。任务服务层专注于细粒度管理，状态流转遵循预设规则防止回退，批量更新采用乐观锁机制避免并发冲突。任务优先级继承自所属项目但允许单独调整，为紧急任务提供灵活调度能力。

\subsubsection{标签系统与多对多关联}

\par 标签系统为用户提供灵活的分类能力，支持跨项目的组织方式。每个标签包含名称和颜色两个属性，颜色以七字符十六进制格式存储，前端直接用于元素样式渲染。标签操作遵循用户隔离原则，所有查询都附加用户标识过滤。项目与标签之间是多对多关系，通过独立的关联表桥接，该关联表冗余存储用户标识以减少权限校验时的表连接操作。标签操作支持幂等性设计，重复分配时静默忽略，批量操作受事务保护。

\par 标签颜色选择遵循色彩无障碍原则，系统提供一组预设颜色覆盖常见视觉区分需求，同时允许用户自定义。标签查询支持分页返回，搜索接口在后端执行模糊匹配，确保高标签数量场景下的流畅使用体验。

\par 从数据模型的角度看，标签与项目之间的多对多关系通过关联表实现，该关联表除了记录项目标识和标签标识外，还冗余存储了用户标识。这种冗余设计使得在权限校验时无需额外的表连接操作，直接通过关联表就能完成用户隔离判断，提升了查询效率。当项目被删除时，级联删除策略自动清理关联表中的对应记录，避免产生悬空关联。

\subsubsection{项目视图设计}

\par 为了优化前端数据获取效率，系统设计了数据库视图聚合项目及其关联数据。该视图通过 JSONB 聚合函数将多行标签和任务记录分别聚合成 JSONB 数组，作为项目的派生列返回。前端通过一次查询即可获得完整信息，避免了传统的 N+1 查询问题。视图底层基表建立了复合索引，覆盖按用户标识加优先级、按用户标识加截止日期等高频查询条件，确保查询响应时间在可接受范围内。

% -----------------------------------------------------------
% §3.6 AI智能体模块设计与实现（核心创新，约4500-6000字）
% -----------------------------------------------------------

\subsection{AI智能体模块设计与实现}
\par 智能体模块是整个系统的技术核心，负责将大语言模型的通用能力转化为面向项目管理场景的实用功能。大语言模型本身是被动的文本生成器，只能根据输入的提示词产生输出，无法主动获取外部信息或执行实际操作。智能体作为大语言模型与外部世界之间的中介，包含感知、决策和执行三个核心组件。本项目选用 Spring AI Alibaba 作为智能体开发框架，该框架提供统一的聊天模型抽象和基于图结构的工作流编排能力，可便捷接入通义千问和 DeepSeek 等大语言模型。

\subsubsection{智能体工作流架构}

\par 智能体工作流采用基于状态图的编排方式。状态图中的节点代表处理步骤，边代表状态转换。工作流从起始节点出发，经过条件判断和处理节点，最终到达结束节点。

\par 工作流的入口是路由节点，该节点承担两项关键的分类决策。第一项决策判断用户当前输入的意图是闲聊还是任务构建。闲聊意图指用户与智能体进行一般性的对话交流，不涉及项目数据的修改；任务构建意图指用户希望智能体帮助创建项目、分解任务或更新项目信息。第二项决策判断当前项目的复杂度级别，如果用户正在讨论一个已有项目，系统分析其描述文本推断复杂度；如果尚未指定项目，则默认按简单项目处理。两项决策的组合产生四种处理分支：简单闲聊、简单构建、复杂闲聊和复杂构建，分别对应四种不同的智能体配置。

\par 路由节点包含自校验机制。在将请求分发到目标智能体之前，轻量级的路由判断模型会评估当前路由决策的合理性。如果模型认为路由不正确，请求会被重新送回路由节点进行二次判断，最多尝试三次后强制进入默认分支。这种设计减少了因路由误判导致的用户体验问题。

\par 除了四种主要的处理智能体外，工作流中还包含一个标题生成智能体。当用户开启新会话时，标题智能体分析用户的初始输入，自动生成简洁的会话标题，帮助用户在会话列表中快速识别不同对话的主题。图 \ref{fig:workflow} 展示了智能体工作流的整体架构。

\begin{figure}[htbp]
    \centering
    \includegraphics[width=\linewidth]{workflow.png}
    \caption{智能体工作流架构图}
    \label{fig:workflow}
\end{figure}

\subsubsection{ReAct智能体实现}

\par ReAct 模式是智能体模块的核心算法范式，其本质是推理与行动交替进行的循环过程。当用户发送一条消息后，智能体首先进入推理阶段，大语言模型分析用户的输入内容、当前的对话历史和可用的工具集合，决定下一步应该采取什么行动。推理的结果被解析为结构化的行动指令，包含行动类型和必要的参数。随后进入行动阶段，系统根据推理结果调用对应的工具函数，将工具的执行结果封装为观察信息。观察信息连同之前的推理过程一起追加到对话上下文中，进入下一轮推理。这个循环持续进行，直到大语言模型判断任务已经完成，此时智能体输出最终的回复文本。

\par 系统提示词在 ReAct 循环中扮演角色定义和行为规范的关键作用，定义了智能体的身份定位、行为边界和输出格式要求。模型参数的配置同样影响着智能体的行为表现：温度参数控制输出文本的随机性，较低的温度使输出更加确定和保守，适合需要精确执行的工具调用场景；较高的温度使输出更加多样和创造性，适合开放式的闲聊对话。系统支持在配置中切换底层模型，当前可选的模型包括通义千问和 DeepSeek，在默认情况下，千问的小模型用于智能路由和上下文压缩，DeepSeek 则负责日常对话和任务拆解。图 \ref{fig:agent-pages} 展示了智能体模块的前端交互页面。

\begin{figure}[htbp]
    \centering
    \includegraphics[width=0.2\linewidth]{create.png}
    \hspace{0.1\linewidth}
    \includegraphics[width=0.2\linewidth]{decomposition.png}
    \caption{智能体模块前端页面}
    \label{fig:agent-pages}
\end{figure}

\par ReAct 模式在实践中面临大语言模型输出不确定性的挑战。系统在提示词中提供了详细的工具使用示例和边界情况处理说明，通过少样本学习引导模型形成稳定的行为模式。同时，系统对模型的输出进行结构化解析，当解析失败时自动重试，而不是直接将原始输出展示给用户。

\subsubsection{情绪调节与认知重构的提示工程实现}

\par 如第2章所述，拖延的本质并非时间管理失败，而是情绪调节失败。因此，智能体的设计不能仅停留在任务拆解和列表管理层面，更需要在交互的每一个触点嵌入情绪安抚与认知重构的能力。本项目通过精细化的提示词工程，将四种差异化的心理干预策略编码到智能体的行为范式中，使智能体在不同任务类型和交互场景下能够自动切换对应的情绪支持模式。

\par 路由节点在实现情绪与认知双维度分流中扮演关键角色。提示词要求模型从信息差与认知边界的视角判断用户任务：当任务的执行路径已锁死、用户比AI更清楚如何完成时（如备考、健身、整理发票），拖延的阻力主要来自情绪而非认知，系统将其标记为简单任务分支，策略是提供极强的情绪鼓励与陪伴感；当任务结构模糊、需要从零到一搭建框架时（如撰写报告、设计页面），拖延的阻力来自认知过载，系统将其标记为复杂任务分支，策略是提供高信息密度的专业拆解。这种基于心理机制的动态分流，确保了后续干预策略与用户真实痛点的精准匹配。

\par 在简单任务分支中，智能体被配置为温暖、随和且极具耐心的行动陪伴者。简单闲聊模式的提示词严格约束智能体的行为边界：禁止使用命令或压迫感词汇，采用口语化、低压的表达方式，每次回复只关注当下这一步，禁止催促或追问为什么要拖延。当用户表现出抗拒时，提示词预设了safety\_net安抚话术，如完全没关系，我们之前的约定就是只要走到门口就可以回来，你已经完成了最难的启动部分。这种设计将五分钟法则和微习惯理论编码为可执行的交互策略，通过极微小的物理动作建议（如站起来、把水杯接满）降低大脑杏仁核的防御反应，从而降低启动阻力。

\par 简单构建模式进一步引入了多巴胺设计和去成败化机制。提示词要求智能体在任务拆解前生成一段极具吸引力的Hook（诱导钩子），从三种策略中选择最适合的一种：科学激励型强调极小投入的巨大回报，用具体数据说服理性大脑；极低门槛型通过去目标化彻底粉碎任务包袱，如我们只是穿上那双舒服的鞋子，下楼呼吸两口新鲜空气；预期快乐型则利用时间旅行技巧，让用户提前感知完成后的爽感。任务拆解遵循物理优先、环境暗示、微小到荒谬的原则，第一步必须是无需动脑的纯物理动作。Safety Net机制明确告知用户只要完成了前两步，你就可以随时停下来，这依然是一次胜利，通过消除对失败的恐惧来降低心理负担。

\par 在复杂任务分支中，智能体被配置为充满好奇心的知识合伙人。复杂闲聊模式的提示词要求智能体扮演小黄鸭调试法中的倾听者角色，允许用户的回答是混乱的、碎片的，通过追问为什么怎么做帮助用户理清思路，而非直接给出完美框架。其工作流程遵循破冰、追问与连接、可视化显现、最小化行动的认知热身路径，逐步降低用户对复杂任务的恐惧感。复杂构建模式则通过三段式心理学梯度实现认知重构：前1至2步安排极低阻力但有具体产出的启动动作，中间3至6步严格套用行业公认的标准工作流，最后1至2步回归交付形式并由AI直接提供防坑点和预测提问。提示词特别约束智能体避免使用第二人称命令式表达，首句客观说明任务内容，后续融入我们营造并肩作战的氛围，以温柔、包容的语气降低用户的心理防御。

\par 前端交互设计对上述情绪支持策略形成了有效配合。如图\ref{fig:agent-pages}所示，AI对话界面采用卡通猫咪头像和温暖的配色方案，营造安全、亲切的对话氛围。任务拆解后的界面以卡片形式呈现每个微步骤，配合进度可视化和显式的正向反馈（如先开始做5分钟按钮），将抽象的情绪调节策略转化为具体的视觉和操作元素，进一步降低了用户的认知负荷和行动门槛。

\subsubsection{工具调用设计}

\par 工具调用是智能体与业务数据交互的唯一通道。系统将工具分为查询类工具和行动类工具两大类。查询类工具仅读取数据，不会对用户的数据产生任何副作用；行动类工具则会修改数据。每个工具函数都通过结构化的描述信息注册到智能体框架中，描述信息包含工具的名称、功能说明、参数列表和每个参数的详细说明。当智能体进入推理阶段时，大语言模型根据这些描述信息判断当前场景下哪个工具最为合适。

\par 查询类工具的设计遵循最小权限原则，智能体在回答用户问题时只获取回答问题所必需的数据，避免将用户的全部项目数据暴露给大语言模型。行动类工具的设计则强调安全性和可预测性，每个修改类工具在执行前都会进行权限校验，确保当前操作用户只能修改自己的数据。工具调用的参数设计遵循最小完备原则，每个工具只接收完成任务所必需的参数，避免参数冗余导致智能体在推理时产生困惑。

\subsubsection{上下文与会话管理}

\par 智能体在处理用户请求时需要访问用户的项目数据以提供上下文感知的回复，但直接将用户的全部数据暴露给大语言模型既不安全也不高效。系统设计了会话级上下文管理机制，用于控制数据的可见范围和缓存策略。

\par 上下文管理器采用线程安全的哈希表实现，以会话标识为键、上下文实例为值。每个活跃会话在首次被访问时自动创建对应的上下文实例，采用懒加载策略，初始为空状态，仅在工具调用触发数据加载时才填充实际内容。上下文的生命周期与会话绑定，当会话被删除或超过一定时间未活跃时，对应的上下文实例从管理器中移除，释放占用的内存资源。

\par 上下文管理还承担着数据新鲜度的维护职责。当用户在对话过程中通过前端界面直接修改了项目或任务信息时，系统在检测到用户数据发生变化时主动使对应会话的上下文缓存失效，强制智能体在下次访问时重新加载最新数据。

\par 大语言模型的对话具有状态依赖性，每一轮回复都依赖于之前的对话历史。在流式响应场景下，用户的网络连接可能在对话中途断开。智能体模块引入了检查点持久化机制，检查点是智能体工作流在运行过程中的状态快照，包含了当前的消息历史、推理过程和工具调用记录。每次智能体完成一轮推理和行动循环后，当前的状态被序列化为字节流存入缓存。当用户重新发起请求时，系统首先检查缓存中是否存在该会话的检查点，如果存在则加载检查点恢复对话状态，智能体从上一次中断的位置继续处理。

\par 会话管理负责维护用户与智能体之间的对话上下文。每个会话对应一次独立的对话主题，包含会话名称和智能体生成的项目草稿。会话表中的草稿字段用于保存智能体在对话过程中为用户构建的项目结构化数据，用户可以在任何时候将草稿转化为正式的项目记录。对话历史的查询采用基于游标的分页方式，使用时间戳作为游标，保证了分页结果的一致性和顺序性。聊天记录的批量操作设计考虑了智能体工作流的特殊性，系统采用批量写入的方式，在智能体完成一轮完整处理后，将所有相关的消息记录打包一次性写入数据库，利用数据库的事务机制保证全部成功或全部回滚。

% -----------------------------------------------------------
% §3.7 安全模块设计与实现
% -----------------------------------------------------------

\subsection{安全模块设计与实现}

\par 安全模块独立于业务模块运行，通过过滤器链机制拦截并校验所有请求。

\subsubsection{JWT过滤器链}

\par 认证过滤器链由三个核心组件构成。第一个是令牌解析过滤器，负责从请求头中提取认证令牌并验证其格式。如果请求头中缺少令牌或格式不正确，过滤器直接返回未认证错误响应，请求不会进入后续业务处理流程。

\par 第二个组件是令牌验证提供者，负责验证令牌的签名和有效期。验证提供者使用与令牌签发时相同的密钥和算法进行签名计算，将结果与令牌携带的签名比对，任何不一致都意味着内容已被篡改。签名验证通过后进一步检查过期时间，确认令牌仍在有效期内。

\par 第三个组件是上下文清理过滤器，在请求处理完成后执行。认证过程中产生的用户身份信息被绑定到当前线程的本地存储，请求结束后清理过滤器将其移除，防止线程复用时身份信息泄漏。过滤器链还配置了白名单机制，登录、注册和健康检查端点无需认证即可访问。白名单采用路径匹配模式，可精确匹配特定路径或批量匹配某一模式下的所有子路径。

\par 过滤器链的执行顺序确保认证逻辑先于业务逻辑。请求首先经过解析和验证，认证失败时在早期就被拦截，不会进入后续资源消耗较大的业务处理环节。这种快速失败机制减少了无效请求的资源占用。

\subsubsection{权限控制}

\par 系统的权限控制分为两个层面。第一个层面是基于角色的访问控制，系统定义了普通用户和管理员两种角色，不同角色可访问的接口范围不同。角色信息作为声明嵌入认证令牌中，验证后提取到当前上下文供后续权限决策使用。

\par 第二个层面是数据级别的权限控制，确保用户只能访问和操作自己的数据。该控制在各业务模块的数据访问层实现，所有涉及用户数据的查询都附加当前登录用户标识作为过滤条件，查询结果只返回属于该用户的数据。即使请求绕过接口层权限校验，数据层的用户隔离条件仍能阻止水平越权访问。

\par 两个层面的权限控制互为补充，形成纵深防御的安全体系。角色层面决定用户能访问哪些接口类型，数据层面决定用户在接口中能看到哪些具体数据。普通用户可访问项目查询接口，但数据隔离确保只能看到自己的项目。管理员虽拥有更广泛的接口权限，访问业务数据时同样受数据层面约束。这种双重防护机制最大程度降低了安全风险。

\par 权限控制的设计还考虑了未来扩展需求。当前系统定义了普通用户和管理员两种角色，新增角色只需在角色枚举中增加条目，无需修改令牌结构和验证逻辑。数据层面的权限控制同样可以扩展，当未来需要支持项目共享功能时，只需在查询条件中增加共享关系的判断，现有的用户隔离机制可作为基础进行扩展。安全模块的整体设计确保了系统在面对各类访问场景时都能维持可控的安全边界。

% -----------------------------------------------------------
% §3.8 语音服务模块设计与实现
% -----------------------------------------------------------

\subsection{语音服务模块设计与实现}

\subsubsection{NLS Token生成服务}

\par 语音服务模块负责为前端提供阿里云自然语言处理服务的访问令牌，支持语音识别和语音合成两项能力。对于注意缺陷多动障碍用户而言，传统的文本输入方式可能存在较高的操作门槛，语音交互能够有效降低这一门槛，使用户可以通过自然的语音方式与系统进行交互，提升系统的易用性和无障碍体验。

\par 由于阿里云自然语言处理服务采用临时令牌机制进行访问控制，客户端无法直接使用长期凭证连接服务，而需要通过后端获取短期有效的访问令牌。令牌的获取流程包含四个步骤：后端使用预配置的访问密钥对调用阿里云安全令牌服务，获取包含凭证信息和权限策略的临时令牌；将临时令牌存入缓存并设置匹配的生存时间；将缓存中的令牌返回给前端；前端使用该令牌直接建立与服务的连接。当令牌接近过期时，前端再次向后端请求，后端优先从缓存返回未过期的令牌。缓存机制避免了每个前端请求都触发阿里云服务调用，显著降低了请求延迟和云服务成本。

\par 这种间接令牌机制具有多重安全优势。长期访问密钥仅保存在后端，不会暴露给客户端，即使前端代码被逆向分析，攻击者也只能获得短期有效的临时令牌。后端在令牌签发时附加细粒度的权限策略，限制令牌只能访问语音识别和合成接口，即使临时令牌泄露，造成的损害也被严格限制在最小范围。此外，缓存机制通过设置合理的生存时间，在保证令牌可用性的同时避免了不必要的阿里云 API 调用，既降低了运营成本，也减少了对外部服务的依赖。

\par 语音服务模块的设计充分考虑了目标用户群体的特殊需求。对于注意缺陷多动障碍用户而言，传统的文本输入方式可能存在较高的操作门槛，长时间的键盘操作容易导致疲劳和注意力分散。语音交互为用户提供了一种更为自然的交互方式，用户可以通过口述的方式创建项目、添加任务和与智能体对话，显著降低了操作复杂度。同时，语音合成能力使得系统能够将文字信息以语音形式播报给用户，在用户不方便查看屏幕的场景下仍能获得信息反馈。

\par 从技术实现角度看，语音服务模块采用了轻量级的服务架构。模块内部仅包含令牌生成服务和缓存管理两个核心组件，令牌生成服务负责调用阿里云安全令牌服务获取临时凭证，缓存管理负责维护令牌的有效期和刷新策略。这种精简的设计使得模块的维护成本极低，同时也为后续扩展预留了空间。当未来需要增加语音流的转发或处理功能时，可以在此模块基础上进行扩展，而不会影响其他业务模块的正常运行。语音服务模块虽然当前仅实现了令牌生成功能，但作为系统多模态交互能力的基础设施，其存在为后续完整的语音交互闭环提供了必要的支撑，是系统面向特殊用户群体进行无障碍设计的重要体现。

% -----------------------------------------------------------
% §3.9 前端设计与接口规划
% -----------------------------------------------------------

\subsection{前端设计与接口规划}

\par 本节阐述前端部分的设计思路、技术选型规划和接口对接方案。前端作为用户与系统交互的直接触点，其设计质量直接影响用户体验和系统的易用性。考虑到目标用户群体在注意力和操作耐心方面可能存在特殊需求，前端设计遵循简洁直观的原则，减少视觉干扰和不必要的操作步骤。

\par 在技术选型方面，前端计划采用现代化的前端框架进行开发，组件化设计使得界面元素可以复用和组合，降低开发维护成本。界面布局以卡片式信息架构为主，每个项目、任务和会话都以独立的卡片呈现，用户可以通过滑动和点击等简单手势完成大部分操作。颜色方案采用低饱和度配色，避免过于鲜艳的色彩刺激，同时保证足够的对比度以确保文字可读性。交互动效保持克制，避免过度动画分散用户注意力。

\par 前端与后端的交互采用两种通信方式。常规的增删改查操作通过标准的 RESTful API 进行，前端发起 HTTP 请求，后端返回 JSON 格式的响应数据。对于 AI 对话场景，由于大语言模型的响应是逐字生成的，AI 对话接口采用服务端推送技术，后端在生成过程中持续向前端发送数据片段，前端即时渲染收到的内容，营造出类似真人打字回复的流式效果，显著提升了用户对系统响应速度的感知。

\par AI 对话界面的设计是整个前端的核心。界面采用经典的聊天式布局，底部为输入区域，上方为消息列表。用户消息和智能体回复分别展示在列表两侧，通过不同的背景色和头像区分发送者身份。当智能体调用工具时，界面显示轻量级加载指示器。工具调用结果以结构化卡片嵌入对话流中，用户可直接点击卡片跳转详情页进行编辑。

\par 项目管理和任务管理界面采用看板式布局。每个项目是一个独立看板，内部按任务状态分多列，用户可拖拽任务卡片在不同状态列之间移动，实现状态快速流转。项目优先级通过卡片边框颜色和排序位置直观呈现，高优先级项目始终排在前列并带有醒目标记。标签以彩色小圆点形式展示在项目卡片上，点击标签可筛选关联项目。

\par 响应式设计确保前端在不同尺寸设备上都能提供可用体验。桌面端采用侧边栏加主内容区布局，充分利用大屏幕空间。移动端将侧边栏折叠为底部导航栏，内容区域全屏展示，手势操作替代鼠标点击。前端设计同时兼顾触摸屏和鼠标键盘两种交互方式，确保无论用户使用何种设备都能顺畅操作。

\par 从接口规划角度看，前端与后端之间通过统一的接口契约进行协作。后端提供的 RESTful API 遵循统一的响应格式，前端可以编写通用的请求处理和错误提示逻辑。接口文档采用 SpringDoc 自动生成，前端开发人员可以在线查看各接口的请求参数、响应结构和示例数据，降低了前后端协作的沟通成本。对于 AI 对话接口，前端需要处理服务端推送的数据流，建立连接、接收数据片段、渲染更新和连接异常处理等逻辑都需要在前端框架中进行妥善设计。由于前端目前尚未进入实际开发阶段，本节所述为设计规划方案，后续开发将以此为基础展开实现。


